\section{Training \& Inference}
\label{sec:training}

Training a streaming 3D reconstruction model end-to-end on long sequences is
challenging: pose errors in early frames propagate along the trajectory and
destabilize the loss landscape, making direct optimization on hundreds of views
impractical.
To address this, we adopt a two-stage training curriculum.
The first stage trains an offline base model on short, diverse multi-view data
to establish robust geometric priors~(\cref{sec:base_training}).
The second stage introduces our proposed Geometric Context Attention (GCA) and
progressively scales to long sequences, transferring the base model's
geometric foundations to the streaming setting~(\cref{sec:streaming_training}).
We describe the training data curation in~\cref{sec:training_data} and
the inference pipeline in~\cref{sec:inference_mode}.


\subsection{Base Model Training}
\label{sec:base_training}

\paragraph{Model Initialization.}
We initialize the ViT backbone from DINOv2~\cite{oquab2023dinov2} 
with a patch size of 14 pixels, followed by 24 alternating blocks of
frame attention and cross-frame attention, following the architecture
of VGGT~\cite{wang2025vggt}.
At this stage, we use standard global attention rather than GCA:
since the training data includes both unordered multi-view collections and
temporally ordered video sequences, global attention imposes no temporal
structure and can fully exploit both data types.
The number of input views per training sample is randomly sampled between
2 and 24, matching the diverse scale of available datasets.

\paragraph{Optimization.}
We use AdamW with a base learning rate of $2 \times 10^{-4}$ and weight
decay of 0.05.
The learning rate follows a linear warmup from $10^{-8}$ to the base rate
over the first 5\% of training, followed by cosine annealing back to
$10^{-8}$ over the remaining 95\%.
Training runs for 160K iterations.

\paragraph{Data Augmentation.}
Images are resized to a maximum dimension of 518 pixels.
To improve robustness to appearance variation across datasets with diverse
capture conditions, we apply aggressive photometric augmentation:
random color jittering (brightness, contrast, saturation $\pm 0.5$;
hue $\pm 0.1$) with probability 0.9, random grayscale conversion with
probability 0.05, and random spatial rescaling in
$[0.8\times, 1.2\times]$ with aspect ratios sampled from $[0.33, 1.0]$.
Additionally, we apply \emph{co-jittering}---applying an identical color
transform to all frames within a scene---with probability 0.3, and
independent per-frame transforms otherwise.
This encourages the model to rely on geometric cues rather than
appearance shortcuts when frames share similar photometric
characteristics, while the independent transforms build robustness to
inter-frame appearance variation.

\paragraph{Distributed Training.}
Training requires approximately 21,500 GPU hours, using fully sharded data
parallelism (FSDP) with gradient checkpointing and bfloat16 mixed precision
to manage memory consumption.


\subsection{Streaming Model Training}
\label{sec:streaming_training}

\paragraph{Initialization from Base Model.}
We initialize the streaming model from the pretrained base model weights
and replace global attention with GCA.
Since the query, key, and value projections in GCA share the same
parameterization as global attention, the pretrained weights transfer
directly, providing a strong initialization that accelerates convergence.

\paragraph{Optimization and Progressive Curriculum.}
We train for 160K iterations with a base learning rate of
$5 \times 10^{-4}$, using the same warmup and cosine annealing schedule
as the first stage.
To stabilize training on progressively longer sequences, we adopt a
view curriculum: the number of training views increases linearly from 24
to 320 over the course of training.
The starting count of 24 views matches the maximum used in the base
stage, and the upper limit of 320 is set by the GPU memory budget under
context parallelism.
Similarly, the local pose reference-window size $k$ of GCA is randomly sampled from 16 to 64
during training, exposing the model to varying receptive fields and
improving robustness at inference time when different window sizes may
be used.

\paragraph{Context Parallelism.}
As the number of views grows, GPU memory becomes the primary bottleneck
due to the quadratic cost of cross-frame attention.
We employ the Ulysses~\cite{jacobs2023deepspeed} context-parallelism
strategy with a parallelism dimension of 16, distributing different views
across GPUs and computing attention via all-to-all collective
communication.
Our implementation builds on TorchTitan~\cite{liang2025torchtitan} and
Magi Attention.
Training requires approximately 15,360 GPU hours.

\subsection{Training Data}
\label{sec:training_data}

\paragraph{Dataset Composition.}
We curate a training corpus of 29 datasets covering indoor, outdoor,
object-centric, synthetic, and real-world scenarios.
The full list, together with data format, scene type, and per-stage
sampling ratios, is given in~\cref{tab:dataset_composition}.
At the coarsest level, the datasets fall into two categories:
\emph{multi-view collections}, where frames are unordered and may lack
temporal continuity, and \emph{video sequences}, where frames follow a
continuous camera trajectory.
This distinction drives the different sampling strategies used in each
training stage.

\paragraph{Stage 1: Diverse Short-Sequence Data.}
The first stage aims to build general geometric priors from a broad
distribution of scenes.
We draw from all 29 datasets with roughly balanced sampling ratios.
The multi-view collections include
BlendedMVS~\cite{yao2020blendedmvs},
HyperSim~\cite{roberts2021hypersim},
MegaDepth~\cite{li2018megadepth},
MVS Synth~\cite{huang2018deepmvs},
GTA SFM~\cite{wang2020flow},
CO3D~\cite{reizenstein2021common},
Objaverse~\cite{deitke2023objaverse},
and Texverse~\cite{zhang2025texverse}.
The video datasets include
Unreal4K~\cite{tosi2021smd},
WildRGBD~\cite{xia2024rgbd},
TartanAir~\cite{wang2020tartanair},
TartanAirV2~\cite{wang2020tartanair},
TartanAirGround~\cite{patel2025tartanground},
PointOdyssey~\cite{zheng2023pointodyssey},
VirtualKITTI~\cite{cabon2020virtual},
Kubric~\cite{greff2022kubric},
DL3DV~\cite{ling2024dl3dv},
Replica~\cite{straub2019replica},
SceneRGBD~\cite{McCormac:etal:ICCV2017},
Mapfree~\cite{arnold2022map},
Aria Synthetic Environments~\cite{pan2023aria},
ADT~\cite{pan2023aria},
ScanNet~\cite{dai2017scannet},
ScanNet++~\cite{yeshwanth2023scannet++},
MatrixCity~\cite{li2023matrixcity},
MidAir~\cite{Fonder2019MidAir},
and our internal game dataset.
Each iteration samples 2 to 24 frames per scene, with a dynamic batch
sampler that packs at most 48 images per GPU.
Frame selection relies on a \emph{nearby sampler}: a reference frame is
chosen at random, and the remaining frames are drawn from a spatial
window around it, without enforcing any temporal order.
This unordered sampling is well-suited to the mixed-modality data in
this stage, where many datasets have no natural frame ordering.

\paragraph{Stage 2: Long-Trajectory Video Data.}
The second stage shifts the distribution toward long, temporally coherent
sequences needed for streaming reconstruction.
We significantly increase the sampling weights for datasets with
extended trajectories and multi-scene coverage, including
TartanAir~\cite{wang2020tartanair},
TartanAirV2,
TartanAirGround~\cite{patel2025tartanground},
MidAir~\cite{Fonder2019MidAir},
MatrixCity~\cite{li2023matrixcity},
Waymo~\cite{sun2020scalability},
VirtualKITTI~\cite{cabon2020virtual},
KITTI-360~\cite{liao2022kitti},
ScanNet++~\cite{yeshwanth2023scannet++},
ScanNet~\cite{dai2017scannet},
and our internal game datasets,
while down-weighting or dropping multi-view-only datasets that lack
temporal structure
(see~\cref{tab:dataset_composition} for exact ratios).

\paragraph{Foldback Video Sampler.}
To produce temporally coherent training subsequences from long videos,
we replace the spatial nearby sampler with a \emph{foldback video
sampler}.
The sampler starts at a random frame and advances with a random stride.
Upon reaching a sequence boundary, it reverses direction and draws a new
stride (distinct from the previous one) to avoid degenerate oscillation.
This mechanism yields subsequences with naturally varying frame rates
and no forward-time bias, providing the model with diverse temporal
contexts during training.

\begin{table}[t]
\centering
\caption{Dataset composition, data format, scene type, and sampling proportions across the two training stages. Blank entries indicate that the dataset was not used in that stage.}
\label{tab:dataset_composition}
\begin{tabular}{llc | cc}
\toprule
Dataset & Format & Scene Type & Stage 1 Ratio (\%) & Stage 2 Ratio (\%) \\
\midrule
BlendedMVS~\cite{yao2020blendedmvs}     & Multi-view & Single & 1.1  & 1.9  \\
HyperSim~\cite{roberts2021hypersim}       & Multi-view & Single & 5.8  & 5.7  \\
MegaDepth~\cite{li2018megadepth}      & Multi-view & Multi  & 3.9  & -    \\
Unreal4K~\cite{tosi2021smd}       & Video      & Single & 1.0  & 0.9  \\
WildRGBD~\cite{xia2024rgbd}       & Video      & Single & 5.8  & 1.9  \\
TartanAir~\cite{wang2020tartanair}      & Video      & Multi  & 3.9  & 7.6  \\
TartanAirV2~\cite{wang2020tartanair}    & Video      & Multi  & 5.8  & 10.8 \\
TartanGround~\cite{patel2025tartanground} & Video     & Multi  & 5.8  & 10.8 \\
Waymo~\cite{sun2020scalability}        & Video      & Multi  & -    & 0.9  \\
PointOdyssey~\cite{zheng2023pointodyssey}   & Video      & Single & 0.3  & 0.9  \\
VirtualKITTI~\cite{cabon2020virtual}  & Video      & Multi  & 0.8  & 1.9  \\
Kubric~\cite{greff2022kubric}         & Video      & Single & 0.3  & 0.9  \\
DL3DV~\cite{ling2024dl3dv}          & Video      & Multi  & 11.0 & 5.7  \\
MVS-Synth~\cite{huang2018deepmvs}      & Multi-view & Single & 1.7  & -    \\
GTA-SFM~\cite{wang2020flow}      & Multi-view & Single & 1.7  & -    \\
CO3D~\cite{reizenstein2021common}          & Multi-view & Single & 5.5  & -    \\
SceneRGBD~\cite{McCormac:etal:ICCV2017}     & Video      & Single & 7.3  & 5.7  \\
Mapfree~\cite{arnold2022map}       & Video      & Multi  & 3.9  & 1.5  \\
Aria Synthetic~\cite{pan2023aria} & Video      & Single & 7.3  & 5.7  \\
ADT~\cite{pan2023aria}            & Video      & Single & 1.0  & -    \\
Objaverse~\cite{deitke2023objaverse}      & Mesh & Single & 2.7  & -    \\
Texverse~\cite{zhang2025texverse}       & Mesh & Single & 2.7  & -    \\
ScanNet++~\cite{yeshwanth2023scannet++}      & Video      & Single & 3.9  & 2.8  \\
ScanNet~\cite{dai2017scannet}        & Video      & Single & 1.9  & 2.8  \\
MatrixCity~\cite{li2023matrixcity}     & Video      & Multi  & 1.7  & 7.6  \\
MidAir~\cite{Fonder2019MidAir}         & Video      & Multi  & 2.9  & 5.7  \\
KITTI-360~\cite{liao2022kitti}      & Video      & Multi  & -    & 3.8  \\
Internal Game  & Video      & Multi  & 10.6 & 10.8 \\
Gibson~\cite{xia2018gibson}  & Mesh      & Multi  & - & 2.6 \\
Matterport3D~\cite{chang2017matterport3d}  & Mesh      & Multi  & - & 2.6 \\
HM3D~\cite{ramakrishnan2021hm3d}  & Mesh      & Multi  & - & 2.6 \\
\bottomrule
\end{tabular}%
\end{table}

\subsection{Inference Modes}
\label{sec:inference_mode}

\method supports two inference modes, Direct Output and Visual
Odometry (VO), that share a common keyframe selection mechanism.

\paragraph{Keyframe Selection.}
When the input sequence exceeds the maximum training length, we
employ an adaptive keyframe selection strategy to control KV-cache
growth.
For each incoming frame, the model first estimates its depth map and
camera pose, then computes the optical flow relative to the most
recent keyframe using the predicted pose and depth.
If the flow magnitude exceeds a predefined threshold, the frame is
designated as a new keyframe, whose features are appended to the
KV cache; otherwise, it is discarded.
This mechanism is shared by both inference modes.

\paragraph{Direct Output Mode.}
Direct mode is the default inference setting.
The model processes frames causally through GCA, with the full
three-level context (anchor, trajectory memory, and local window)
accumulating continuously without reset.
Each frame directly outputs an absolute camera pose and a dense
depth map.
In this mode, prediction errors accumulate solely from the model's
frame-by-frame inference, without introducing any additional error
from external alignment steps.
Although the model is trained on sequences of up to 320 views, we
empirically find that the Direct mode with keyframe selection remains
stable for approximately $10\times$ the training length
(${\sim}3{,}000$ frames), beyond which prediction quality gradually
degrades.

\paragraph{Visual Odometry (VO) Mode.}
For sequences that far exceed the Direct mode's effective range
(e.g., tens of thousands of frames), we switch to VO mode.
The input is partitioned into overlapping local windows.
Within each window, an initial subset of frames is processed jointly
to establish a robust local scale and coordinate system, and the
remaining frames are processed causally via GCA with keyframe
selection.
At the end of each window, the model state is reset.
To fuse successive windows into a single global trajectory, we
compute a Sim(3) alignment between the overlapping regions of
consecutive windows, recovering the relative scale, rotation, and
translation.
This enables \method to process arbitrarily long sequences with
bounded memory, at the cost of additional drift introduced at each
window boundary: unlike Direct mode, VO mode incurs extra alignment
error that compounds with the number of windows.

\paragraph{Trade-offs.}
Direct mode produces more accurate trajectories by avoiding
inter-window alignment error, and is the preferred choice when the
sequence length stays within ${\sim}3{,}000$ frames.
For inputs that substantially exceed this range, VO mode generalizes
more effectively at the cost of accumulated alignment drift at the 
window boundaries.
In practice, the choice is determined by the sequence length and the
required level of global consistency.

\paragraph{Default Inference Configuration.}
Unless otherwise specified, all experiments in this report use Direct
Output Mode with a local pose-reference window size $k=64$ and
keyframe interval $m=1$, at a resolution of $518 \times 378$ with
bfloat16 precision.
For the large-scale demo videos that span city-scale environments
or extremely long sequences, we use VO mode.
